%==========================================================================
%  TRES0312 Fysikk — Foiler-mal
%
%  MØrk/lys modus: kommenter ut én av de to linjene nedenfor.
%==========================================================================
\newif\ifdarkmode
\darkmodetrue      % ← mørk modus
%\darkmodefalse   % ← lys modus  (fjern % foran og legg % foran linjen over)

\documentclass[aspectratio=169, 12pt]{beamer}

\usepackage{amd-beamer}
\usetikzlibrary{overlay-beamer-styles}
\usetikzlibrary{calc}

%--------------------------------------------------------------------------
%  DOKUMENTINFO
%--------------------------------------------------------------------------
\title{Oppdrift og Arkimedes' lov}
\subtitle{TRES0312 --- kapittel 7.4}
\author{Ben David Normann}
\date{\today}

%--------------------------------------------------------------------------
%  SEKSJONSKONTEKST
%--------------------------------------------------------------------------
\setcounter{section}{7}   % Kapittel 7: "Fluiddynamikk"

\begin{document}

%--------------------------------------------------------------------------
%  SLIDE 1 — TITTELSIDE
%--------------------------------------------------------------------------
\begin{frame}[plain]
  \titlepage
  \dekkerdelkap{kapittel 7.4}
  \vspace{0.3cm}
  \begin{center}
  {\bfseries\color{repcol} Første time: oppdriftskraft. Etter pausen: flyter eller synker?}
  \end{center}
\end{frame}

%--------------------------------------------------------------------------
%  SLIDE 2 — Tre kjappe fra forrige gang
%--------------------------------------------------------------------------
\begin{frame}{}

  \repetisjon{Tre kjappe fra forrige gang}{%
    \begin{enumerate}[1.]
      \item Hva sier Pascals lov?
      \item I en hydraulisk jekk er $A_2$ ti ganger så stort som $A_1$. Hvor mye større blir $F_2$ enn $F_1$?
      \item Får vi noe «gratis» med en hydraulisk jekk?
    \end{enumerate}
  }

\end{frame}

%==========================================================================
%  DEL 1: OPPDRIFTSKRAFT
%==========================================================================

%--------------------------------------------------------------------------
%  SLIDE 3 — Til ettertanke: steiner under vann
%--------------------------------------------------------------------------
\begin{frame}{Oppdrift}

  \tanke{}{%
  Hvorfor føles en stein lettere når du løfter den under vann?
  }
  \pause
  \svar{
  Trykket i vannet øker med dybden. Derfor presser vannet hardere opp på undersiden av steinen enn ned på oversiden. Resultatet er en kraft oppover: \textit{oppdriftskraften}.
  }

\end{frame}

%--------------------------------------------------------------------------
%  SLIDE 4 — Definisjon: oppdriftskraft
%--------------------------------------------------------------------------
\begin{frame}{Oppdrift}

  \defnat{7.4.1}{Oppdriftskraft}{
  Oppdriftskraften på et legeme i en væske i ro er
  \[
    O = \rho_v\,\mathrm{g}\,V_\text{f}
  \]
  der $\rho_v$ er tettheten til væsken og $V_\text{f}$ er volumet av væsken som legemet fortrenger.
  }

  \pause
  \merk{Fortrengt volum}{%
  $V_\text{f}$ er den delen av legemet som er under væskeoverflaten. Er legemet helt under, er $V_\text{f}$ lik hele volumet.
  }

\end{frame}

%--------------------------------------------------------------------------
%  SLIDE 5 — Aktivitet: stein i vann
%--------------------------------------------------------------------------
%  Fasit: G = 6.0*9.81 = 58.9 N; O = 1000*9.81*2.5e-3 = 24.5 N ≈ 25 N;
%         S = G - O = 34.3 N ≈ 34 N (i lufta: 59 N).
\begin{frame}{Aktivitet}

  \aktivitet{Stein i vann}{%
  Du holder en stein i en snor, slik at den henger i ro helt under vann. Steinen har masse $6.0\,\mathrm{kg}$ og volum $2.5\,\mathrm{L}$.
  \begin{enumerate}[a)]
    \item Tegn kreftene som virker på steinen.
    \item Finn oppdriftskraften.
    \item Finn snorkraften. Sammenlign med snorkraften i lufta.
  \end{enumerate}
  }

\end{frame}

%--------------------------------------------------------------------------
%  SLIDE 6 — Oppsummering: oppdriftskraft
%--------------------------------------------------------------------------
\begin{frame}{}

  \oppsum{Oppdriftskraft}{%
  \begin{enumerate}[1.]
    \item[] Oppdriften skyldes at trykket er større under legemet enn over.
    \item[] $O = \rho_v\,\mathrm{g}\,V_\text{f}$, der $V_\text{f}$ er volumet under væskeoverflaten.
  \end{enumerate}
  }

\end{frame}

%--------------------------------------------------------------------------
%  SLIDE 7 — PAUSE
%--------------------------------------------------------------------------
\begin{frame}
\begin{center}
  \Huge{Pause}\\[8pt]
  {\large Etter pausen: flyter eller synker?}
\end{center}
\end{frame}

%==========================================================================
%  DEL 2: FLYTER ELLER SYNKER? ARKIMEDES' LOV
%==========================================================================

%--------------------------------------------------------------------------
%  SLIDE 8 — Til ettertanke: stålskip
%--------------------------------------------------------------------------
\begin{frame}{Flyter eller synker?}

  \tanke{}{%
  Stål har nesten åtte ganger så stor tetthet som vann. Hvordan kan et skip av stål flyte?
  }
  \pause
  \svar{
  Det er skipet \textit{som helhet} som teller: stål pluss all lufta inni. Den gjennomsnittlige tettheten er mindre enn tettheten til vann.
  }

\end{frame}

%--------------------------------------------------------------------------
%  SLIDE 9 — Flyter eller synker?
%--------------------------------------------------------------------------
\begin{frame}{Flyter eller synker?}

  For et legeme helt under væsken er $G = \rho_L\,\mathrm{g}\,V$ og $O = \rho_v\,\mathrm{g}\,V$:
  \graabox{$\displaystyle F_\text{netto} = O - G = (\rho_v - \rho_L)\,\mathrm{g}\,V$}
  \pause
  \begin{itemize}
    \item $\rho_L < \rho_v$: legemet \textbf{flyter}.
    \item $\rho_L = \rho_v$: legemet \textbf{svever}.
    \item $\rho_L > \rho_v$: legemet \textbf{synker}.
  \end{itemize}
  \pause
  Når et legeme flyter i ro, er $O = G$.

\end{frame}

%--------------------------------------------------------------------------
%  SLIDE 10 — Aktivitet: isbit i et glass vann
%--------------------------------------------------------------------------
%  Fasit: O = G gir V_f/V = rho_is/rho_vann = 0.92/1.0 = 0.92 = 92 %.
%         Bonus: isfjell i sjøvann: 0.92/1.03 = 0.89 = 89 %.
\begin{frame}{Aktivitet}

  \aktivitet{Isbit i et glass vann}{%
  En isbit ($\rho_\text{is} = 920\,\mathrm{kg/m^3}$) flyter i et glass vann ($\rho_\text{vann} = 1.0\cdot 10^3\,\mathrm{kg/m^3}$).
  \begin{enumerate}[a)]
    \item Hvor stor andel av isbiten ligger under vann? \textit{(Hint: sett $O = G$.)}
    \item Et isfjell flyter i sjøvann ($1.03\cdot 10^3\,\mathrm{kg/m^3}$). Hvor stor andel ligger under vann?
  \end{enumerate}
  }

\end{frame}

%--------------------------------------------------------------------------
%  SLIDE 11 — Definisjon: Arkimedes' lov
%--------------------------------------------------------------------------
\begin{frame}{Arkimedes' lov}

  \defnat{7.4.2}{Arkimedes' lov}{
  Oppdriftskraften på et legeme i en væske i ro er lik tyngden av den fortrengte væsken:
  \[
    O = G_\text{f} = \rho_v\,\mathrm{g}\,V_\text{f}
  \]
  }

\end{frame}

%--------------------------------------------------------------------------
%  SLIDE 12 — Aktivitet: oppdrift på en båt
%--------------------------------------------------------------------------
%  Fasit: V_f = 0.15*12 = 1.8 m^3; O = 1000*9.81*1.8 = 1.77e4 N ≈ 1.8e4 N;
%         m = O/g = 1.8e3 kg.
\begin{frame}{Aktivitet}

  \aktivitet{Oppdrift på en båt}{%
  En båt har et volum på $12\,\mathrm{m^3}$, og $15\,\%$ av volumet ligger under vann.
  \begin{enumerate}[a)]
    \item Finn oppdriftskraften på båten.
    \item Hva er massen til båten?
  \end{enumerate}
  }

\end{frame}

%--------------------------------------------------------------------------
%  SLIDE 13 — Oppsummering: flyter eller synker
%--------------------------------------------------------------------------
\begin{frame}{}

  \oppsum{Flyter eller synker?}{%
  \begin{enumerate}[1.]
    \item[] Sammenlign tettheten til legemet med tettheten til væsken.
    \item[] Et legeme som flyter i ro, har $O = G$.
    \item[] Arkimedes' lov: oppdriften er lik tyngden av den fortrengte væsken.
  \end{enumerate}
  }

\end{frame}

%--------------------------------------------------------------------------
%  SLIDE 14 — Neste gang
%--------------------------------------------------------------------------
\begin{frame}
\begin{center}
  \Huge{Neste gang: termofysikk}
  \vspace{6mm}
\end{center}
\end{frame}

\end{document}
